
\documentclass{article}
\usepackage{colortbl}
\usepackage{makecell}
\usepackage{multirow}
\usepackage{supertabular}

\begin{document}

\newcounter{utterance}

\centering \large Interaction Transcript for game `cladder', experiment `full\_v1.5\_default', episode 9577 with reflexion{-}Qwen3.5{-}4B.
\vspace{24pt}

{ \footnotesize  \setcounter{utterance}{1}
\setlength{\tabcolsep}{0pt}
\begin{supertabular}{c@{$\;$}|p{.15\linewidth}@{}p{.15\linewidth}p{.15\linewidth}p{.15\linewidth}p{.15\linewidth}p{.15\linewidth}}
    \# & \multicolumn{2}{c}{Player} && \multicolumn{2}{c}{Game Master} \\
    \hline

    \theutterance \stepcounter{utterance}  
    & & & \multicolumn{4}{p{0.6\linewidth}}{
        \cellcolor[rgb]{0.9,0.9,0.9}{
            \makecell[{{p{\linewidth}}}]{
                \texttt{\tiny{[P1$\langle$GM]}}
                \texttt{You have attempted this same task before and did not succeed. These are the plans you wrote after those attempts.} \\
\texttt{Trial \#1: New plan: I failed to complete the task because I cut off my reasoning mid{-}sentence while analyzing the counterfactual mechanism for "Tijv." The prompt requires a complete answer starting with "Yes" or "No" followed by full reasoning. My strategy was to construct the causal graph correctly but I abandoned the analysis of the specific counterfactual condition (\$\textbackslash{}neg \textbackslash{}text\{Xevo\}\$) before deriving the final conclusion about "Gyzp." Next time, I will ensure I finish the entire logical deduction chain: explicitly state the status of Tijv given \$\textbackslash{}neg \textbackslash{}text\{Xevo\}\$, determine if Gyzp is still caused (via the direct link from Tijw or the indirect link from Tijv), and provide a definitive "Yes" or "No" based on whether the absence of Xevo (and thus Tijv) combined with the absence of Tijw guarantees the absence of Gyzp. I must avoid truncating the explanation and ensure the final verdict is clearly stated.} \\
\\ 
\\ 
\texttt{You are an expert in causal inference. The following question is not a typical commonsense query, but rather a meticulously designed question created by a professor specializing in causal inference, intended to assess the students' mastery of the course content.} \\
\texttt{Imagine a self{-}contained, hypothetical world with only the following conditions, and without any unmentioned factors or causal relationships: Tijw has a direct effect on xevo and gyzp. Xevo has a direct effect on tijv. Tijv has a direct effect on gyzp. Tijw is unobserved. We know that tijw causes xevo. xevo causes tijv. tijw or tijv causes gyzp. We observed an individual is not tijw. Would an individual is not gyzp if not xevo instead of xevo?} \\
\texttt{Start your answer with 'Yes' or 'No', followed by additional reasoning or evidence to support your explanation.} \\
            }
        }
    }
     \\ \\

    \theutterance \stepcounter{utterance}  
    & \multicolumn{4}{p{0.6\linewidth}}{
        \cellcolor[rgb]{0.561,0.737,0.561}{
            \makecell[{{p{\linewidth}}}]{
                \texttt{\tiny{[P1$\rangle$GM]}}
                \texttt{Yes.} \\
\\ 
\texttt{**Reasoning:**} \\
\\ 
\texttt{1.  **Analyze the Causal Structure:**} \\
\texttt{    *   **Tijw** is an unobserved root node.} \\
\texttt{    *   **Xevo** is a direct child of **Tijw** (\$Tijw \textbackslash{}to Xevo\$).} \\
\texttt{    *   **Tijv** is a direct child of **Xevo** (\$Xevo \textbackslash{}to Tijv\$).} \\
\texttt{    *   **Gyzp** has two parents: **Tijw** and **Tijv** (\$Tijw \textbackslash{}to Gyzp\$ and \$Tijv \textbackslash{}to Gyzp\$). The prompt states "Tijw or tijv causes gyzp," which in the context of causal graphs with multiple parents implies that Gyzp is caused if *either* parent is active (logical OR).} \\
\\ 
\texttt{2.  **Analyze the Observed Condition:**} \\
\texttt{    *   We observe an individual who is **not Tijw** (\$\textbackslash{}neg Tijw\$).} \\
\texttt{    *   The question asks: "Would an individual be **not Gyzp** if **not Xevo** instead of **Xevo**?"} \\
\texttt{    *   This phrasing asks us to evaluate the counterfactual outcome: What is the status of Gyzp in a world where \$Xevo\$ is false (\$\textbackslash{}neg Xevo\$), given that we already know \$Tijw\$ is false?} \\
\\ 
\texttt{3.  **Evaluate the Counterfactual Scenario (\$\textbackslash{}neg Xevo\$):**} \\
\texttt{    *   **Effect on Xevo:** In this hypothetical scenario, \$Xevo\$ is false (\$\textbackslash{}neg Xevo\$).} \\
\texttt{    *   **Effect on Tijv:** Since \$Xevo\$ causes \$Tijv\$ (\$Xevo \textbackslash{}to Tijv\$), if \$Xevo\$ is false, then \$Tijv\$ must also be false (assuming no other causes for Tijv, which the prompt confirms by stating "without any unmentioned factors"). Thus, \$\textbackslash{}neg Xevo \textbackslash{}implies \textbackslash{}neg Tijv\$.} \\
\texttt{    *   **Effect on Gyzp:** The prompt states that Gyzp is caused by \$Tijw\$ OR \$Tijv\$.} \\
\texttt{        *   We know \$Tijw\$ is false (from the observation).} \\
\texttt{        *   In this counterfactual scenario, we derived that \$Tijv\$ is false.} \\
\texttt{        *   Since both potential causes of Gyzp (\$Tijw\$ and \$Tijv\$) are false, the condition for Gyzp to occur is not met.} \\
\\ 
\texttt{4.  **Conclusion:**} \\
\texttt{    *   With \$\textbackslash{}neg Tijw\$ and \$\textbackslash{}neg Xevo\$ (which necessitates \$\textbackslash{}neg Tijv\$), there are no active causes for Gyzp.} \\
\texttt{    *   Therefore, Gyzp will not occur.} \\
            }
        }
    }
    & & \\ \\

    \theutterance \stepcounter{utterance}  
    & & & \multicolumn{2}{p{0.3\linewidth}}{
        \cellcolor[rgb]{0.9,0.9,0.9}{
            \makecell[{{p{\linewidth}}}]{
                \texttt{\tiny{[GM$|$GM]}}
                \texttt{Yes} \\
            }
        }
    }
    & & \\ \\

    \theutterance \stepcounter{utterance}  
    & & & \multicolumn{2}{p{0.3\linewidth}}{
        \cellcolor[rgb]{0.9,0.9,0.9}{
            \makecell[{{p{\linewidth}}}]{
                \texttt{\tiny{[GM$|$GM]}}
                \texttt{yes} \\
            }
        }
    }
    & & \\ \\

    \theutterance \stepcounter{utterance}  
    & & & \multicolumn{2}{p{0.3\linewidth}}{
        \cellcolor[rgb]{0.9,0.9,0.9}{
            \makecell[{{p{\linewidth}}}]{
                \texttt{\tiny{[GM$|$GM]}}
                \texttt{game\_result = WIN} \\
            }
        }
    }
    & & \\ \\

\end{supertabular}
}

\end{document}
